\section{Lean proofs and checked answers}
\label{sec:certification}

Agreement between reasoners is the usual evidence that a classification is
correct, and it fails when two reasoners share a defect.  KM therefore
connects its answers to a Lean~4 development \cite{demoura2021lean4} in two
ways: theorems about the reasoning rules, and checkers that examine the
evidence for one answer.  The development contains no admitted proofs, and
its public theorems depend only on Lean's standard axioms
\texttt{propext}, \texttt{Classical.choice}, and \texttt{Quot.sound}.

The development defines a clause language with a first-order interpretation
semantics, the EL completion rules, hypertableau evidence, derivations of the
context calculus, and the OWL~2 functional syntax with its direct semantics
for the constructs of the fragments discussed below.  For each production
procedure it defines an evidence format and an executable checker, and it
proves that an accepted evidence document implies the answer.  A positive
answer (a subsumption, a type, or an inconsistency) is accepted with a
derivation whose every step is an instance of a rule proved sound.  A negative
answer (a non-subsumption, a non-type, or consistency) is accepted with a
finite model of the input in which the answer is visible: KM produces the
model from its result rows or from its model export, and the checker verifies
that the model satisfies the input.  Checking a model requires no proof that
the procedure which found it is complete.  For an inconsistent input the
evidence is a refutation.

Completeness is proved for reference procedures on fragments.  A reference
procedure is a direct implementation of a calculus inside Lean, with proofs of
soundness and completeness with respect to the OWL~2 direct semantics; it is
exponential and answers small inputs.  The development contains reference
procedures for OWL~EL, $\mathcal{ALC}$, and $\mathcal{SHI}$, and v1.7.0
extends the completeness results to $\mathcal{SHIQ}$ at the level of normal
forms, to knowledge bases with an ABox, to datatypes, and to a fragment with
nominals.  \todo{State each fragment and theorem precisely from
\texttt{docs/V1.7.0-LEAN-INVENTORY.md} and \texttt{docs/CERTIFICATION-ROADMAP.md}.}
The optimised Rust procedures carry no completeness proof; outside the
fragments above, and for inputs too large for a reference procedure, the
evidence for a negative answer is the checked model.

The checkers are optional at run time and are enabled by one setting.  With
the checkers enabled, a run publishes an answer together with the checker
verdicts, and an answer whose evidence a checker rejects is withheld.  A
checker can also decline an input or exhaust its time or memory budget, in
which case the answer is published as solved and uncertified.
Section~\ref{sec:eval-cert} reports how many answers on the ORE pool are
certified.  The benchmark timings of Section~\ref{sec:evaluation} are measured
with the checkers disabled, which is the default configuration.

The trusted base of a certified answer consists of Lean's kernel and compiler,
the statement of the OWL~2 semantics in Lean, and the code that reads the
source ontology into the checker.  The Rust procedures, the router, and the
Konclude port are outside the trusted base of a certified answer, because the
checker examines their output.  For an uncertified answer the evidence is the
agreement with other reasoners and the differential tests reported in
Section~\ref{sec:eval-diff}.  Five gate scripts build the development and the
checkers from the tagged source, fail on any admitted proof, list the axioms of
every public theorem, run fixtures in which a valid evidence document is
altered and must be rejected, and run the Rust integration tests that pass
worker evidence through the checkers.  Every release passes these gates on the
tagged commit; for v1.7.0 the gates comprise 658 tests, alongside 3,818 Rust
tests and 219 Java tests.
